\chapter{Compter sans compter}

Ce chapitre ne contient aucune formule. Il contient l'idée entière.

\section{La flaque dans la cour}

Imaginez une flaque d'eau de forme bizarre au milieu d'une cour rectangulaire.
Vous voulez savoir quelle \textbf{surface} elle occupe. Vous n'avez pas de
règle, et la flaque ne ressemble à aucune figure connue : ce n'est ni un carré,
ni un cercle, ni rien qui ait une formule.

Alors vous faites ceci. Vous prenez une poignée de cailloux et vous les lancez
\textbf{au hasard, partout dans la cour}, sans viser. Puis vous comptez ceux qui
ont fait « plouf ».

\begin{nkretenir}[Le raisonnement complet]
Si 30 cailloux sur 100 sont tombés dans la flaque, alors la flaque occupe
environ \textbf{30 \% de la cour}. Et la cour, elle, est un rectangle : sa
surface, vous la connaissez.

\[ \text{surface de la flaque} \;\approx\; 0{,}30 \times \text{surface de la cour} \]
\end{nkretenir}

Vous venez de mesurer l'immesurable sans jamais approcher la flaque. C'est
\emph{tout} Monte-Carlo. Le nom vient du casino de Monte-Carlo : on joue aux
dés pour obtenir un résultat sérieux.

Deux choses sont déjà vraies, et elles ne changeront plus de tout le cours :

\begin{itemize}
  \item avec 10 cailloux, la réponse est mauvaise ; avec 10\,000, elle est
        bonne. \textbf{Plus on lance, plus c'est juste} ;
  \item on n'aura \textbf{jamais} la réponse exacte. On aura une réponse
        presque juste, dont l'erreur rétrécit tant qu'on continue.
\end{itemize}

\section{Le premier vrai calcul : trouver $\pi$ avec de la pluie}

Passons à un cas où la réponse est connue, pour vérifier que la méthode ne
ment pas.

Prenez un carré de côté 2, allant de $-1$ à $+1$ dans les deux directions, et
le cercle de rayon 1 qui tient exactement dedans.

\begin{itemize}
  \item surface du carré $= 2 \times 2 = 4$ ;
  \item surface du cercle $= \pi \times r^2 = \pi$ (puisque $r = 1$).
\end{itemize}

Un point lancé au hasard dans le carré a donc une chance sur $4/\pi$ de tomber
dans le cercle — autrement dit, la proportion de points dans le cercle vaut
$\pi/4$. Retournez la relation :

\[ \pi \;\approx\; 4 \times \frac{\text{points tombés dans le cercle}}{\text{points lancés}} \]

Reste à savoir si un point $(x, y)$ est dans le cercle. Il l'est si sa distance
au centre ne dépasse pas 1, c'est-à-dire si $x^2 + y^2 \leq 1$. On ne prend
même pas la racine carrée : comparer les carrés donne le même verdict et coûte
moins cher.

\begin{nkcode}{Estimer $\pi$ — le programme complet}
float32 EstimerPi(nk_uint32 n) {
    nk_uint32 dedans = 0;
    for (nk_uint32 i = 0; i < n; ++i) {
        const float32 x = math::NkRand.NextFloat32(-1.f, 1.f);
        const float32 y = math::NkRand.NextFloat32(-1.f, 1.f);
        // Pas de racine carree : comparer les carres donne le meme verdict.
        if (x * x + y * y <= 1.f)
            ++dedans;
    }
    return 4.f * (float32)dedans / (float32)n;
}
\end{nkcode}

Voici ce que rend ce programme, en pratique :

\begin{center}
\begin{tabular}{@{}rll@{}}
\toprule
\textbf{Tirages} & \textbf{Estimation} & \textbf{Erreur} \\
\midrule
100          & 3,08     & 2 \% \\
10\,000      & 3,132    & 0,3 \% \\
1\,000\,000  & 3,1416   & 0,02 \% \\
\bottomrule
\end{tabular}
\end{center}

\begin{nkretenir}[Ce que vous venez de faire]
Vous avez calculé $\pi$ \textbf{sans jamais écrire $\pi$}, et sans connaître la
moindre formule de géométrie du cercle. Vous n'avez utilisé qu'une chose : la
capacité à répondre à la question « ce point est-il dedans ? ».

C'est le motif universel. Monte-Carlo ne demande jamais de savoir résoudre le
problème — seulement de savoir \textbf{évaluer une réponse ponctuelle}.
\end{nkretenir}

\section{Pourquoi ça marche vraiment}

Ce n'est pas une coïncidence ni un bricolage. C'est la \emph{loi des grands
nombres} : quand on répète une expérience au hasard un grand nombre de fois, la
moyenne des résultats se rapproche de la valeur vraie.

Dit autrement : chaque caillou pris isolément ne vous apprend presque rien
(il est dedans ou dehors, c'est tout). Mais leur \textbf{moyenne} sait quelque
chose qu'aucun caillou ne sait.

\begin{nkpiege}[Ce n'est pas « ça finit par tomber juste par chance »]
Une confusion fréquente : croire que le résultat oscille et qu'on attrape la
bonne valeur au passage. Non — la moyenne \textbf{converge}. Chaque nouveau
tirage tire l'estimation vers la vraie valeur, et l'écart rétrécit de façon
prévisible (chapitre 3). Ce n'est pas de la chance, c'est une garantie
mathématique.
\end{nkpiege}

\begin{nkexo}[Vérifier la convergence de ses propres yeux]
Modifiez \texttt{EstimerPi} pour qu'elle affiche l'estimation courante tous les
1000 tirages, sur 100\,000 tirages. Tracez la courbe.

Vous verrez trois choses : elle démarre n'importe où, elle se resserre vite au
début, puis elle rampe. Ce ralentissement est le sujet du chapitre 3, et c'est
le prix à payer de la méthode.
\end{nkexo}
